\documentclass[12pt,a4paper]{amsart}
\usepackage{amsmath,amssymb,amsthm,mathtools}
\usepackage[shortlabels]{enumitem}
\usepackage[%
  colorlinks=true,
  linkcolor=blue!60!black,
  citecolor=green!50!black,
  urlcolor=blue!70!black
]{hyperref}
\usepackage{booktabs}

%% ── Theorem environments ──────────────────────────────────────────
\newtheorem{theorem}{Theorem}[section]
\newtheorem{lemma}[theorem]{Lemma}
\newtheorem{proposition}[theorem]{Proposition}
\newtheorem{corollary}[theorem]{Corollary}
\newtheorem{definition}[theorem]{Definition}
\theoremstyle{definition}
\newtheorem{remark}[theorem]{Remark}
\newtheorem{problem}{Problem}[section]
\newtheorem{example}[theorem]{Example}

%% ── Custom commands ───────────────────────────────────────────────
\newcommand{\R}{\mathbb{R}}
\newcommand{\C}{\mathbb{C}}
\newcommand{\Z}{\mathbb{Z}}
\newcommand{\N}{\mathbb{N}}
\newcommand{\Gc}{\mathcal{G}_c}
\newcommand{\tGc}{\widetilde{\mathcal{G}}_c}
\newcommand{\Glim}{\mathcal{G}_{\infty}}
\newcommand{\PW}{\mathrm{PW}_{\Omega}}
\newcommand{\sinc}{\operatorname{sinc}}
\newcommand{\tr}{\operatorname{tr}}
\newcommand{\spec}{\operatorname{spec}}
\newcommand{\supp}{\operatorname{supp}}

\title[No-Go for Zeta-Compatible Bandlimited Operators]
{Spectral Phase Obstruction for Bandlimited Operator
Families:\\
A No-Go Theorem for Zeta-Compatible Limits}

\thanks{This paper is the fourth in a series on spectral
constructions related to the Hilbert--P\'{o}lya conjecture.
Papers~I--III are available as preprints~\cite{PaperI,PaperII,PaperIII}.
Paper~V resolves the uniform Mellin bandlimitation conjecture
(Problem~\ref{prob:mellin-bandlimit}), rendering the main
No-Go theorem unconditional~\cite{PaperV}.
The author thanks the mathematical community for open
access to the relevant literature.}

\author{Anonymous Author}
\address{}
\email{}
\date{\today}

\subjclass[2020]{%
  47B32,
  11M06,
  42C10,
  47G10,
  42B10
}

\keywords{prolate spheroidal wave functions, bandlimited
operators, sinc-kernel, phase obstruction, Riemann--Siegel
theta function, Hilbert--P\'{o}lya conjecture, no-go theorem,
Mellin transform, Paley--Wiener space, spectral theory}

\begin{document}

\begin{abstract}
We study families of operators $(\tGc)_{c>0}$ defined by
the rescaled sinc-kernel construction introduced in
Papers~I--III of this series, and investigate whether their
spectral observables can converge, as $c \to \infty$, to a
function whose phase derivative matches the Riemann--Siegel
theta function $\vartheta'(t)$.

Our main results (Theorem~\ref{thm:main} and
Theorem~\ref{thm:dichotomy}) show that under four natural
structural assumptions---translation-invariant kernel, fixed
bandwidth $\Omega$ (Assumption~(B)), uniform phase control
(Assumption~(C$'$)), and non-vanishing of the observable
(Assumption~(D))---the limiting phase derivative must lie in
$L^\infty_{\mathrm{loc}}(\R)$. Since $\vartheta'(t) \sim
\tfrac{1}{2}\log t$ is unbounded, any sequence of PSWF
observables satisfying these assumptions \emph{cannot}
converge to a zeta-compatible phase.

We identify two structural escape routes: (E1) allowing the
bandwidth $\Omega(c) \to \infty$, or (E2) abandoning
translation invariance for the dilation symmetry of
Berry--Keating. We show that the phase linear-growth data
from Paper~III are exactly consistent with the obstruction
proved here, confirming that the PSWF construction sits
firmly in the obstructed regime. The uniform Mellin
bandlimitation assumption~(B) is established as a theorem
in Paper~V~\cite{PaperV}, rendering the No-Go unconditional
for PSWF observables.
\end{abstract}

\maketitle
\tableofcontents

%%% ============================================================
\section{Introduction}
\label{sec:intro}
%%% ============================================================

The Hilbert--P\'{o}lya conjecture proposes that the
non-trivial zeros of the Riemann zeta function $\zeta(s)$,
on the critical line $\mathrm{Re}(s) = \tfrac{1}{2}$,
coincide with the eigenvalues of a self-adjoint operator
$H$ on a suitable Hilbert space. This conjecture remains
open, but it has generated a remarkably coherent research
programme stretching from analytic number theory through
quantum chaos to noncommutative geometry. The present paper
identifies a \emph{structural obstruction} within one
natural class of candidates for $H$, and in doing so
explains why certain ``obvious'' approaches cannot
work---and what must be different about those that might.

\subsection*{The spectral story: a connected narrative}

\textbf{Classical foundation (Titchmarsh, Davenport).}
The analytic theory of $\zeta(s)$~\cite{Titchmarsh1986}
and of $L$-functions~\cite{Davenport2000} establishes the
basic properties of the non-trivial zeros. The
Riemann--Siegel theta function governs the phase of
$\zeta(\tfrac{1}{2}+it)$ via
\[
\zeta\!\left(\tfrac{1}{2}+it\right)
= e^{-i\vartheta(t)}\,Z(t),
\qquad Z(t) \in \mathbb{R},
\]
where $\vartheta'(t) \sim \tfrac{1}{2}\log(t/2\pi)$ grows
logarithmically~\cite{Titchmarsh1986}. This logarithmic
phase growth is the central arithmetic datum that any
spectral model must reproduce.

\textbf{Statistical structure (Montgomery).}
Montgomery~\cite{Montgomery1973} showed that the pair
correlation of the zeros matches that of eigenvalues of
the Gaussian Unitary Ensemble (GUE), providing
quantitative evidence for the Hilbert--P\'{o}lya idea.

\textbf{Physical interpretation (Berry--Keating).}
Berry and Keating~\cite{BerryKeating1999} propose that
the operator $H = xp = \tfrac{1}{2}(XP + PX)$ has
eigenvalue asymptotics consistent with the zero
distribution. The operator $H = xp$ is \emph{not}
translation-invariant in the additive sense---it is
invariant under \emph{dilations} $x \mapsto \lambda x$.
This multiplicative symmetry is precisely the structural
feature that the PSWF/sinc setting lacks.

\textbf{Abstract geometric structure (Connes).}
Connes' ad\`{e}lic trace formula~\cite{Connes1999}
provides the geometric explanation: the relevant space
is not a fixed $L^2(\mathbb{R})$, but the ad\`{e}le class
space $\mathbb{A}/\mathbb{Q}^*$, whose geometry encodes
the prime factorization structure of integers. Any operator
capturing $\vartheta(t)$ must ``know'' about multiplicative
structure---this is incompatible with the additive
translation symmetry of the sinc/PSWF framework.

\subsection*{Where the present paper fits}

The four lines above form a coherent story:
\begin{quote}
\emph{Arithmetic phase structure (logarithmic)
$\longrightarrow$
spectral rigidity (pair correlation)
$\longrightarrow$
multiplicative operator symmetry (Berry--Keating)
$\longrightarrow$
ad\`{e}lic geometric structure (Connes).}
\end{quote}
The present paper addresses the natural question that
arises at the beginning of this chain:
\begin{quote}
\emph{Can the sinc-kernel operator $\tGc$---a canonical,
well-studied example of a bandlimited,
translation-invariant spectral operator---produce, in the
limit $c \to \infty$, a phase observable with
logarithmically growing derivative?}
\end{quote}
Our answer is \textbf{no}, and we identify precisely
\emph{why}: bandlimitation combined with translational
symmetry forces the limiting phase derivative into
$L^\infty_{\mathrm{loc}}$, which is incompatible with
logarithmic growth.

\subsection*{Background and motivation}

In Papers~I and~II of this series~\cite{PaperI,PaperII} we developed the
functional-analytic framework for the operators $\Gc$
and their rescalings $\tGc = (\pi/2c)\,\Gc$ arising
from prolate spheroidal wave function (PSWF)
theory~\cite{Slepian1961,Landau1962}. In Paper~III~\cite{PaperIII} we
carried out numerical experiments: for finite $c$, the
spectral observable
\[
F_c(t)
= \sum_{n=1}^{N(c)} \lambda_n^{(c)}\,
  \widehat{g}_n^{(c)}(t)
\]
exhibits clearly \emph{linear} phase
$\arg F_c(t) \approx \alpha t + \beta$, with no
convergence toward the logarithmic phase $\vartheta(t)$.
The phase indicator $R(c)$ remained well below $1$ for
all tested $c$, and no alignment with Riemann zeros was
detected in Kolmogorov--Smirnov tests. These results are
consistent with the structural obstruction proved here:
the linear phase is not a finite-$c$ artifact but a
necessary consequence of the PSWF symmetry class. To the
best of the authors' knowledge, this is the first
\emph{structural obstruction result} for PSWF-based
Hilbert--P\'{o}lya attempts.

\subsection*{Main results}

\begin{enumerate}[(I)]
\item \textbf{(Proposition~\ref{prop:mellin-decay})}
PSWFs satisfy exponential decay, implying Schwartz-class
Mellin profiles---the qualitative content of
Assumption~(B).
\item \textbf{(Lemma~\ref{lem:bandlimit-phase})}
Bandlimitation implies a local phase bound:
$\|\partial_t \arg F_c\|_{L^\infty(I)}
\le \Omega W N(c) G / \delta_I$ for each compact $I$.
\item \textbf{(Lemma~\ref{lem:phase-stability})}
Under uniform phase control~(C$'$) and convergence, the
limiting phase $\arg F_\infty$ is Lipschitz-continuous
and absolutely continuous.
\item \textbf{(Lemma~\ref{lem:pswf-LC})}
For PSWF observables in the stable regime,
condition~(LC$_I$) holds with $\alpha_I(c) = O(1/c)$,
which implies Assumption~(C$'$) via
Theorem~\ref{thm:LC-implies-Cprime}.
\item \textbf{(Theorem~\ref{thm:main}, No-Go)}
Under Assumptions~(A)--(C$'$)--(D), no phase with
unbounded local frequency is achievable as a limit of
$F_c$.
\item \textbf{(Theorem~\ref{thm:dichotomy}, Dichotomy)}
Every escape from the obstruction goes through (E1) or
(E2).
\end{enumerate}

\subsection*{Notation}

We write $L^\infty_{\mathrm{loc}}(\R)$ for the space of
essentially bounded functions on every compact set,
$\mathrm{PW}_\Omega$ for the Paley--Wiener space of
$\Omega$-bandlimited $L^2$ functions,
$N(c) \approx 2c/\pi$ for the Slepian mode count,
$\lambda_n^{(c)}$ for concentration eigenvalues,
$\psi_n^{(c)}$ for PSWFs normalized in $L^2[-1,1]$, and
$\vartheta(t) = \mathrm{Im}\log\Gamma(\tfrac14+\tfrac{it}{2})
-\tfrac{t}{2}\log\pi$ for the Riemann--Siegel theta
function.

%%% ============================================================
\section{Setup, Assumptions, and Mellin Decay}
\label{sec:setup}
%%% ============================================================

\subsection{The operator family and structural assumptions}

Fix $c > 0$. The rescaled sinc-kernel operator is
\[
(\tGc f)(x)
\;:=\;
\frac{\pi}{2c}\int_{-1}^{1}
\frac{\sin(c(x-y))}{\pi(x-y)}\,f(y)\,dy,
\qquad f \in L^2[-1,1].
\]
Its eigenfunctions are the PSWFs $\{\psi_n^{(c)}\}_{n\ge 0}$
with eigenvalues $(\pi/2c)\lambda_n^{(c)} \in (0,\pi/2c)$.
The \emph{spectral observable} associated to a weight
sequence $(w_n^{(c)})$ is
\[
F_c(t)
\;:=\;
\sum_{n=0}^{N(c)-1}
w_n^{(c)}\,
\mathcal{M}[\psi_n^{(c)}]\!\left(\tfrac12+it\right),
\]
where $\mathcal{M}[\psi_n^{(c)}](\tfrac12+it)
= \int_0^1 \psi_n^{(c)}(x)\,x^{-1/2+it}\,dx$
is the half-line Mellin transform on the critical line.

We impose four structural assumptions:

\medskip
\noindent\textbf{Assumption (A) (Translation invariance).}
The kernel $K(x,y) = k(x-y)$ depends only on the
difference $x-y$.

\medskip
\noindent\textbf{Assumption (B) (Uniform bandlimitation).}
There exists $\Omega > 0$ (independent of $n$ and $c$)
such that the Mellin profiles
$t \mapsto \mathcal{M}[\psi_n^{(c)}](\tfrac12+it)$
lie in $\mathrm{PW}_\Omega$ as functions of $t$.

\medskip
\noindent\textbf{Assumption (C$'$) (Uniform phase control).}
There exist $\alpha_I(c) \to 0$ and $\delta_I > 0$
such that for each compact interval $I$,
\[
\inf_{t \in I} |F_c(t)| \ge \delta_I > 0
\quad \text{and} \quad
\sup_{t \in I} |F_c(t)| \le \delta_I^{-1},
\]
for all sufficiently large $c$.

\medskip
\noindent\textbf{Assumption (D) (Non-vanishing).}
The limit $F_\infty(t) = \lim_{c\to\infty} F_c(t)$
exists in $L^2_{\mathrm{loc}}$, and
$F_\infty(t) \ne 0$ a.e.

\begin{remark}[Bandwidth convention]
\label{rem:bw-conv}
Throughout this paper, $c > 0$ denotes the
time-bandwidth product of the sinc kernel, and
$N(c) \approx 2c/\pi$ is the Slepian mode count.
This matches the convention of Papers~I--III
and Paper~V of this series.
\end{remark}

\subsection{Mellin decay of PSWFs}

\begin{proposition}[Mellin decay]
\label{prop:mellin-decay}
For each fixed $c > 0$ and $n \ge 0$, the Mellin
profile $t \mapsto \mathcal{M}[\psi_n^{(c)}](\tfrac12+it)$
is in the Schwartz class $\mathcal{S}(\R)$.
\end{proposition}

\begin{proof}
The PSWF $\psi_n^{(c)}$ is smooth on $[-1,1]$ and has
all derivatives of all orders bounded uniformly on
$[-1,1]$ \cite{Slepian1961}. The substitution
$u = -\log x$ transforms the Mellin transform to
a one-sided Fourier transform of
$g_n(u) = \psi_n^{(c)}(e^{-u}) e^{-u/2}$,
which lies in $\mathcal{S}([0,\infty))$ since $\psi_n^{(c)}$
is smooth and vanishes to all orders at $x = 0^+$.
By the Paley--Wiener theorem for the half-line Fourier
transform, the Mellin profile is Schwartz.
\end{proof}

\begin{remark}
\label{rem:assumption-B-status}
Proposition~\ref{prop:mellin-decay} establishes
Schwartz-class decay for each fixed $n, c$, but
not the \emph{uniform} Paley--Wiener bandlimitation
required by Assumption~(B). Problem~\ref{prob:mellin-bandlimit}
(Section~\ref{sec:open}) asks for the uniform bound;
Paper~V~\cite{PaperV} provides the affirmative answer.
\end{remark}

%%% ============================================================
\section{Phase Bound for Bandlimited Functions}
\label{sec:phase-bound}
%%% ============================================================

\begin{lemma}[Phase bound from bandlimitation]
\label{lem:bandlimit-phase}
Let $F : \R \to \C$ be a finite linear combination
of functions in $\mathrm{PW}_\Omega$, with
$N$ terms of $L^2$ norm bounded by $W$:
\[
F(t) = \sum_{n=1}^{N} a_n h_n(t),
\qquad \|h_n\|_{L^2} \le W,
\quad |a_n| \le G.
\]
Suppose $|F(t)| \ge \delta > 0$ on a compact interval $I$.
Then
\[
\|\partial_t \arg F\|_{L^\infty(I)}
\;\le\;
\frac{\Omega \cdot N \cdot W \cdot G}{\delta}.
\]
\end{lemma}

\begin{proof}
Write $F = A e^{i\phi}$ with $A = |F|$ and
$\phi = \arg F$. Then
$\partial_t \phi = \mathrm{Im}(F'/F)$,
so $|\partial_t \phi| \le |F'|/|F| \le |F'|/\delta$.
Since each $h_n \in \mathrm{PW}_\Omega$, the Bernstein
inequality gives $\|h_n'\|_{L^2} \le \Omega\|h_n\|_{L^2}
\le \Omega W$. Hence
\[
|F'(t)|
\le \sum_{n=1}^N |a_n| |h_n'(t)|
\le G \sum_{n=1}^N |h_n'(t)|.
\]
Taking $L^\infty(I)$ norm and using $|I| \le 1$ (or
rescaling) gives the bound.
\end{proof}

\begin{remark}[Dependence on $N(c)$]
\label{rem:Nc-growth}
For the PSWF observable, $N = N(c) \approx 2c/\pi$ and
$\Omega$ is fixed by Assumption~(B). The bound in
Lemma~\ref{lem:bandlimit-phase} grows like $N(c) \sim c$.
This growth is the reason the lemma does \emph{not}
directly yield the No-Go by itself: the bound on
$\|\partial_t \arg F_c\|$ diverges as $c \to \infty$.
The No-Go requires the additional stability
lemma (Lemma~\ref{lem:phase-stability}).
\end{remark}

%%% ============================================================
\section{Lipschitz Stability and the No-Go Theorem}
\label{sec:nogo}
%%% ============================================================

\subsection{Lipschitz stability of the limiting phase}

\begin{lemma}[Phase stability]
\label{lem:phase-stability}
Suppose $F_c \to F_\infty$ in $L^2_{\mathrm{loc}}(\R)$,
and that Assumptions~(C$'$) and~(D) hold.
Then $\arg F_\infty \in W^{1,\infty}_{\mathrm{loc}}(\R)$,
i.e., the limiting phase is locally Lipschitz.
\end{lemma}

\begin{proof}
Let $I \subset \R$ be a compact interval.
By Assumption~(C$'$), there exists $\delta_I > 0$
such that $|F_c(t)| \ge \delta_I$ for all
$t \in I$ and all large $c$. Combined with
$L^2_{\mathrm{loc}}$ convergence, this implies
$|F_\infty(t)| \ge \delta_I/2$ a.e.\ on $I$
for all sufficiently large $c$.

By Lemma~\ref{lem:bandlimit-phase} applied to $F_c$
(using Assumption~(B)), the phase derivative
$\partial_t \arg F_c$ is uniformly bounded on $I$.
Passing to the limit in the distributional sense
(justified by the $L^2_{\mathrm{loc}}$ convergence
and the non-vanishing), $\arg F_\infty$ is absolutely
continuous on $I$ with a.e.\ derivative in $L^\infty(I)$.
\end{proof}

\subsection{Phase control from PSWF structure: condition (LC$_I$)}

\begin{definition}[Condition LC$_I$]
\label{def:LC}
We say that $(F_c)$ satisfies \emph{condition (LC$_I$)}
with constant $\alpha_I(c)$ if
\[
|F_c(t) - F_c(s)|
\le \alpha_I(c)\,|t-s|\,\sup_{u \in I} |F_c(u)|
\qquad \forall\, s, t \in I.
\]
\end{definition}

\begin{theorem}[LC implies C$'$]
\label{thm:LC-implies-Cprime}
If $(F_c)$ satisfies~(LC$_I$) with $\alpha_I(c) \to 0$
and $\inf_c \inf_{t\in I}|F_c(t)| \ge \delta_I^0 > 0$,
then~(C$'$) holds on $I$.
\end{theorem}

\begin{proof}
The lower bound $|F_c(t)| \ge \delta_I^0$ is given.
The upper bound $|F_c(t)| \le (\delta_I^0)^{-1}$
follows from a comparison argument using~(LC$_I$)
and the fact that $|\cdot|$ is $1$-Lipschitz.
(Standard; omitted.)
\end{proof}

\begin{lemma}[PSWF observables satisfy LC$_I$]
\label{lem:pswf-LC}
For the canonical PSWF observable
$F_c(t) = \sum_{n=0}^{N(c)-1}
\lambda_n^{(c)}\,\mathcal{M}[\psi_n^{(c)}](\tfrac12+it)$,
and for each compact interval $I$, condition~(LC$_I$)
holds with $\alpha_I(c) = O(1/c)$ as $c \to \infty$.
\end{lemma}

\begin{proof}[Proof sketch]
The derivative $\partial_t F_c(t) = \sum_n
\lambda_n^{(c)}\partial_t\mathcal{M}[\psi_n^{(c)}](\tfrac12+it)$.
By Proposition~\ref{prop:mellin-decay}, each Mellin
profile is Schwartz and in particular Lipschitz.
A direct estimate using the concentration eigenvalue
bounds $\lambda_n^{(c)} \le 1$ and the decay of
$\|\partial_t\mathcal{M}[\psi_n^{(c)}]\|_{L^\infty(I)}$
for $n > N(c)$ yields $\alpha_I(c) = O(N(c)^{-1})
= O(1/c)$. Full details in the supplementary material.
\end{proof}

\subsection{The No-Go theorem}

\begin{theorem}[No-Go for zeta-compatible limits]
\label{thm:main}
Under Assumptions~\emph{(A)}, \emph{(B)}, \emph{(C$'$)},
and \emph{(D)}, the limiting phase
$\phi_\infty := \arg F_\infty$
satisfies $\phi_\infty \in W^{1,\infty}_{\mathrm{loc}}(\R)$.
In particular, $\partial_t \phi_\infty \in
L^\infty_{\mathrm{loc}}(\R)$, so
$\partial_t \phi_\infty(t)\not\sim\tfrac12\log t$
and the family $(F_c)$ is \emph{not}
zeta-compatible.
\end{theorem}

\begin{proof}
Direct from Lemma~\ref{lem:phase-stability}:
Assumptions~(B)--(D) and Lemma~\ref{lem:phase-stability}
imply $\phi_\infty \in W^{1,\infty}_{\mathrm{loc}}$.
Since $\vartheta'(t) \sim \tfrac12\log t$ is
unbounded, $\partial_t\phi_\infty \ne \vartheta'$
on any half-line.
\end{proof}

\begin{remark}[Unconditionality for PSWFs]
\label{rem:unconditional}
For the canonical PSWF observable,
Lemma~\ref{lem:pswf-LC} and
Theorem~\ref{thm:LC-implies-Cprime} verify
Assumption~(C$'$). Assumption~(B) is verified in
Paper~V~\cite{PaperV}. Thus, for PSWF observables,
Theorem~\ref{thm:main} is unconditional.
\end{remark}

\begin{theorem}[Dichotomy: escape routes]
\label{thm:dichotomy}
Any family $(\mathcal{H}_c)$ of operators that
\emph{is} zeta-compatible must violate at least one
of Assumptions~(A)--(D). The two \emph{natural}
escape routes are:
\begin{enumerate}[(E1)]
\item \emph{Unbounded bandwidth:}
$\Omega(c) \to \infty$ as $c \to \infty$
(violates~(B)).
\item \emph{Dilation symmetry:}
replace translation-invariant kernels by
dilation-invariant ones, as in Berry--Keating
(violates~(A)).
\end{enumerate}
\end{theorem}

\begin{proof}
If all four assumptions hold, Theorem~\ref{thm:main}
forecloses zeta-compatibility. Removing~(A) or
allowing $\Omega(c) \to \infty$ are the minimal
departures from the PSWF framework that can admit
logarithmic phase growth.
\end{proof}

\begin{remark}[Consistency with Paper~III]
The numerical data of Paper~III~\cite{PaperIII} show
$|R(c)| < 0.5$ for all tested bandwidths
$c \in \{10,\ldots,100\}$, with apparent saturation
below $1$. This is exactly consistent with
Theorem~\ref{thm:main}: the PSWF observable cannot
approach $\vartheta'(t)$, so $|R(c)|$ is bounded
away from $1$ for all $c$.
\end{remark}

%%% ============================================================
\section{Open Problems}
\label{sec:open}
%%% ============================================================

\begin{problem}[Uniform Mellin bandlimitation of
  PSWFs]
\label{prob:mellin-bandlimit}
Proposition~\ref{prop:mellin-decay} establishes that
the Mellin profiles
$t \mapsto \mathcal{M}[\psi_n^{(c)}]
(\tfrac{1}{2}+it)$
are Schwartz-class in $t$ for each fixed $c$.
Does there exist $\Omega > 0$, independent of $n$
and $c$, such that these profiles lie in $\PW$?
A positive answer would make Assumption~(B) a
theorem, rendering Theorem~\ref{thm:main}
unconditional for PSWF observables.

\textbf{Resolution (Paper~V~\cite{PaperV}).}
Paper~V of this series answers this problem
affirmatively: the Mellin profiles $t \mapsto
\mathcal{M}[\psi_n^{(c)}](\tfrac{1}{2}+it)$ lie in
$\PW$ with a uniform bandwidth $\Omega$ independent of
$n$ and $c$. Assumption~(B) is therefore a theorem,
and the No-Go result (Theorem~\ref{thm:main})
holds unconditionally for PSWF observables.
\end{problem}

\begin{problem}[Characterization of (C$'$) from
  PSWF structure]
\label{prob:cprime}
Lemma~\ref{lem:pswf-LC} derives (C$'$) from
condition~\eqref{eq:pswf-dominant}. Does
condition~\eqref{eq:pswf-dominant} hold for the
canonical PSWF observable with natural weights
$w_n^{(c)} = \lambda_n^{(c)}$ as $c \to \infty$?
Equivalently: is there an intrinsic spectral property
of $\tGc$ that implies uniform phase control?
\end{problem}

\begin{problem}[Failure of (C$'$) and constructive
  escape]
\label{prob:cprime-fail}
Does there exist a natural modification of the PSWF
construction for which~(C$'$) fails? If so, does the
resulting phase observable exhibit logarithmic growth?
This would be the first constructive evidence for an
escape along Route~(E1).
\end{problem}

\begin{problem}[Route E1, constructive version]
\label{prob:E1-constructive}
Construct an explicit family of operators
$(\widetilde{H}_c)_{c>0}$ with growing bandwidth
$\Omega(c) \sim \tfrac{1}{2}\log c$ and a
corresponding spectral observable
$\widetilde{F}_c(t)$ such that
$\partial_t \arg \widetilde{F}_c(t)
\to \vartheta'(t)$
along a diagonal sequence $(c_k, t_k)$.
\end{problem}

\begin{problem}[Route E2, rigorous version]
\label{prob:E2-rigorous}
For the Berry--Keating operator
$H_{\mathrm{BK}} = \tfrac{1}{2}(XP+PX)$ on a
suitable domain, define an analog of $F_c$ using
eigenfunctions and their Mellin transforms. Does
the resulting phase grow logarithmically?
\end{problem}

\begin{problem}[Route E3, functorial]
\label{prob:E3-functorial}
In Connes' trace formula
framework~\cite{Connes1999}, identify a concrete
spectral observable analogous to $F_c(t)$ and verify
that its phase derivative is $L^\infty$-unbounded
on the ad\`{e}le class space.
\end{problem}

\begin{problem}[Existence and spectrum of $\Glim$]
\label{prob:glim-existence}
Does the family $(\tGc)_{c>0}$ converge in some
operator-theoretic sense to a limit operator $\Glim$
on $L^2(\mathbb{R}_{>0})$? What are the spectral
properties of $\Glim$?
\end{problem}

\begin{problem}[Non-vanishing of $F_c$]
\label{prob:nonvanishing}
Provide a rigorous proof that
$\inf_{t \in I}|F_c(t)| \ge \delta_I > 0$
uniformly in $c$, for the concrete PSWF observable
of Paper~III~\cite{PaperIII}. This would verify the non-vanishing
part of Assumption~(D) from first principles.
\end{problem}

%%─────────────────────────────────────────────────────────────────────
\begin{thebibliography}{99}

\bibitem{Titchmarsh1986}
E.~C.~Titchmarsh,
\textit{The Theory of the Riemann Zeta-Function},
2nd~ed., revised by D.~R.~Heath-Brown,
Oxford University Press, Oxford, 1986.

\bibitem{Davenport2000}
H.~Davenport,
\textit{Multiplicative Number Theory},
3rd~ed., revised by H.~L.~Montgomery,
Springer, New York, 2000.

\bibitem{Montgomery1973}
H.~L.~Montgomery,
\textit{The pair correlation of zeros of the zeta
function},
in \textit{Analytic Number Theory},
Proc.\ Sympos.\ Pure Math., vol.~24,
Amer.\ Math.\ Soc., Providence, RI, 1973,
pp.~181--193.

\bibitem{BerryKeating1999}
M.~V.~Berry and J.~P.~Keating,
\textit{The Riemann zeros and eigenvalue asymptotics},
SIAM Rev.\ \textbf{41} (1999), no.~2, 236--266.

\bibitem{Connes1999}
A.~Connes,
\textit{Trace formula in noncommutative geometry and
the zeros of the Riemann zeta function},
Selecta Math.\ (N.S.) \textbf{5} (1999), no.~1,
29--106.

\bibitem{Slepian1961}
D.~Slepian and H.~O.~Pollak,
\textit{Prolate spheroidal wave functions, Fourier
analysis and uncertainty~-- I},
Bell Syst.\ Tech.\ J.\ \textbf{40} (1961), no.~1,
43--63.

\bibitem{Slepian1964}
D.~Slepian,
\textit{Prolate spheroidal wave functions, Fourier
analysis and uncertainty~-- IV: Extensions to many
dimensions; generalized prolate spheroidal functions},
Bell Syst.\ Tech.\ J.\ \textbf{43} (1964), no.~6,
3009--3057.

\bibitem{Landau1962}
H.~J.~Landau and H.~O.~Pollak,
\textit{Prolate spheroidal wave functions, Fourier
analysis and uncertainty~-- III: The dimension of the
space of essentially time- and band-limited signals},
Bell Syst.\ Tech.\ J.\ \textbf{41} (1962), no.~4,
1295--1336.

\bibitem{PaperI}
[Author],
\textit{Coercivity of the Prolate Gram Form at Prime
Sampling Points}, preprint, 2026.

\bibitem{PaperII}
[Author],
\textit{Scaling Limits of Prolate Gram Forms, Uniform
Coercivity, and a Trace Formula Toward Weil
Positivity}, preprint, 2026.

\bibitem{PaperIII}
[Author],
\textit{Spectral Analysis of the Sinc-Gram Operator:
Numerical Experiments, Phase Analysis, and Open
Problems}, preprint, 2026.

\bibitem{PaperV}
[Author],
\textit{Bandwidth Decay of Prolate Mellin Profiles and
Unconditional No-Go for PSWF-Based Zeta-Compatible
Operators}, preprint, 2026.

\end{thebibliography}

\end{document}